\documentclass{article}
\usepackage[utf8]{inputenc}
\usepackage{amsmath,amssymb,amsfonts}
\usepackage{graphicx}
\usepackage{hyperref}
\usepackage{booktabs}
\usepackage[margin=1in]{geometry}
\usepackage{enumitem}
\usepackage{float}
\usepackage{xcolor}
\usepackage{colortbl}
\usepackage{natbib}

% Define colors for matrix notation
\definecolor{gradcolor}{RGB}{76, 175, 80}   % Green for gradient
\definecolor{curlcolor}{RGB}{244, 67, 54}   % Red for curl
\definecolor{harmcolor}{RGB}{33, 150, 243}  % Blue for harmonic

\title{The Shape of Good Behavior: A Geometric Approach to AI Alignment}
\author{Michael Lee}
\date{\today}

\begin{document}

\maketitle

\begin{abstract}
Reinforcement Learning from Human Feedback (RLHF) has become the dominant paradigm for aligning large language models with human preferences. However, the standard approach---learning a scalar reward model from pairwise comparisons---fundamentally assumes that preferences are \emph{transitive} and \emph{consistent}. In practice, human preferences exhibit cyclic inconsistencies (Condorcet paradoxes), context-dependence, and evaluator disagreement, leading to reward hacking and unstable learning dynamics.

We introduce \textbf{Sheaf-Theoretic Reward Spaces (STRS)}, a mathematical framework that models the \emph{preference aggregation problem} as a topological structure before applying geometric optimization. Our key insight is that existing methods---PPO's clipped trust regions and CPO's cost constraints---are \emph{implicit approximations} to navigating a reward manifold with topological defects. We make this structure explicit using Hodge decomposition to separate consistent preferences (gradient) from cyclic inconsistencies (curl), and sheaf cohomology to measure global obstructions to reward aggregation.

We propose \textbf{Sheaf-Geodesic Policy Optimization (SGPO)}, which models safety constraints as geometric singularities creating hard guarantees via infinite geodesic distance. Preliminary experiments on synthetic benchmarks show SGPO detects preference cycles with 71\% accuracy (vs.\ PPO's implicit exploitation) and reduces catastrophic safety violations by 79\% compared to unconstrained baselines.
\end{abstract}

\section{The Preference Aggregation Problem}

Before introducing our geometric framework, we identify the fundamental \emph{computation object} that RLHF attempts to construct: a consistent global preference ordering from noisy, partial, and potentially contradictory local comparisons.

\subsection{What Existing Methods Implicitly Compute}

\textbf{Proximal Policy Optimization (PPO)} \citep{schulman2017proximal} constrains policy updates to a trust region, clipping the probability ratio to $[1-\epsilon, 1+\epsilon]$. This clipping acts as a \emph{defect-limited approximation}: by restricting updates to local neighborhoods, PPO implicitly assumes the reward function is consistent within small regions, avoiding the need to reconcile globally inconsistent preferences. When the reward landscape contains Condorcet cycles, PPO exploits them blindly---accumulating infinite perceived reward while making no real progress.

\textbf{Constrained Policy Optimization (CPO)} \citep{achiam2017constrained} adds cost constraints $\mathbb{E}[C] \le d$ via Lagrangian relaxation. This represents a \emph{constraint-limited approximation}: safety is enforced only for constraints the trainer explicitly specifies. Catastrophic states not anticipated by the constraint set remain reachable with non-zero probability.

\textbf{Our Contribution}: We make the implicit geometric structure explicit. The preference aggregation problem has a natural topological formulation where:
\begin{enumerate}[leftmargin=*]
    \item \textbf{Consistent preferences} form a \emph{gradient field} (derivable from a scalar potential)
    \item \textbf{Cyclic preferences} form a \emph{curl field} (local rotational inconsistencies)
    \item \textbf{Global obstructions} form a \emph{harmonic field} (irreducible topological defects)
\end{enumerate}

\subsection{Topological Features as Safety Diagnostics}

Social choice theory \citep{condorcet1785} tells us that aggregating preferences often leads to paradoxes. We interpret these paradoxes as \emph{topological features} encoded in the Betti numbers ($H_n$) of the reward manifold:

\begin{itemize}[leftmargin=*]
    \item \textbf{$H_0$ (Connected Components):} Count of disconnected ``islands'' of aligned behavior. $H_0 > 1$ indicates fragmented alignment with no safe path between optima.
    
    \item \textbf{$H_1$ (Cyclic Inconsistencies):} Non-vanishing $H^1 \neq 0$ captures Condorcet cycles ($A \succ B \succ C \succ A$). These are not noise---they reveal genuine value conflicts or context-dependence in the preference data.
    
    \item \textbf{$H_2$ (Voids/Black Holes):} Cavities surrounded by high-reward regions but containing catastrophic outcomes---``event horizons'' that are easy to enter but impossible to escape.
\end{itemize}

\begin{figure}[H]
    \centering
    \includegraphics[width=0.8\textwidth]{../../fig1_sheaf_structure.png}
    \caption{Sheaf structure over the reward manifold. Local sections (human feedback patches) must satisfy compatibility conditions to glue into global sections (consistent value functions). Non-trivial $H^1$ indicates sections that cannot be glued.}
    \label{fig:sheaf_structure}
\end{figure}

\section{Hodge Decomposition for Preference Analysis}

We now formalize how to decompose preference feedback into consistent and inconsistent components using combinatorial Hodge theory \citep{jiang2011statistical}.

\subsection{The Preference Graph}

Given pairwise comparisons $(i, j, y_{ij})$ where $y_{ij} > 0$ means item $j$ is preferred over $i$, we construct:
\begin{itemize}[leftmargin=*]
    \item \textbf{Vertices} $V$: Items being compared (states, actions, or trajectories)
    \item \textbf{Edges} $E$: Pairs with observed comparisons
    \item \textbf{Triangles} $T$: Triplets where all three pairwise comparisons exist
\end{itemize}

\subsection{Matrix Formulation with Color-Coded Components}

The preference flow $\mathbf{Y} \in \mathbb{R}^{|E|}$ decomposes via the Hodge-Helmholtz theorem:

\begin{equation}
\mathbf{Y} = \underbrace{\textcolor{gradcolor}{\mathbf{D}_0 \mathbf{s}}}_{\text{\textcolor{gradcolor}{Gradient}}} + \underbrace{\textcolor{curlcolor}{\mathbf{D}_1^\top \mathbf{v}}}_{\text{\textcolor{curlcolor}{Curl}}} + \underbrace{\textcolor{harmcolor}{\mathbf{h}}}_{\text{\textcolor{harmcolor}{Harmonic}}}
\end{equation}

Where:
\begin{itemize}[leftmargin=*]
    \item $\textcolor{gradcolor}{\mathbf{D}_0} \in \mathbb{R}^{|E| \times |V|}$: Boundary operator (edges $\to$ vertices)
    \item $\textcolor{gradcolor}{\mathbf{s}} \in \mathbb{R}^{|V|}$: Scalar potential (\emph{consistent value function})
    \item $\textcolor{curlcolor}{\mathbf{D}_1} \in \mathbb{R}^{|T| \times |E|}$: Boundary operator (triangles $\to$ edges)
    \item $\textcolor{curlcolor}{\mathbf{v}} \in \mathbb{R}^{|T|}$: Triangle potential (\emph{local rotations})
    \item $\textcolor{harmcolor}{\mathbf{h}} \in \ker(\mathbf{D}_0^\top) \cap \ker(\mathbf{D}_1)$: Harmonic flow (\emph{global cycles})
\end{itemize}

\textbf{Key Insight}: The $\textcolor{gradcolor}{\text{gradient component}}$ is what standard RLHF tries to learn. The $\textcolor{curlcolor}{\text{curl component}}$ represents local inconsistencies that can be averaged out. The $\textcolor{harmcolor}{\text{harmonic component}}$ represents \emph{irreducible} global cycles---these cannot be removed without additional context or constraints.

\subsection{Computing the Decomposition}

The decomposition is computed via the graph Laplacians:
\begin{align}
\textcolor{gradcolor}{\mathbf{s}} &= (\mathbf{D}_0^\top \mathbf{D}_0)^\dagger \mathbf{D}_0^\top \mathbf{Y} & \text{(Least-squares potential)} \\
\textcolor{curlcolor}{\mathbf{v}} &= (\mathbf{D}_1 \mathbf{D}_1^\top)^\dagger \mathbf{D}_1 (\mathbf{Y} - \mathbf{D}_0 \mathbf{s}) & \text{(Residual curl)} \\
\textcolor{harmcolor}{\mathbf{h}} &= \mathbf{Y} - \mathbf{D}_0 \mathbf{s} - \mathbf{D}_1^\top \mathbf{v} & \text{(Remaining harmonic)}
\end{align}

The magnitude $\|\textcolor{harmcolor}{\mathbf{h}}\|$ directly measures the $H^1$ cohomology---the degree to which preferences are \emph{fundamentally} inconsistent.

\subsection{Unifying Diverse Feedback Types}

A critical question: how do we construct the preference graph from different feedback types? The Hodge decomposition operates on edge flows, but real RLHF systems receive diverse signals:

\begin{itemize}[leftmargin=*]
    \item \textbf{Pairwise comparison} ($A > B$): Direct edge with $Y_{AB} = 1$
    \item \textbf{K-wise ranking} ($A > B > C$): $\binom{K}{2}$ edges from all pairs
    \item \textbf{Binary feedback} (thumbs up/down): Node weight $r_i$, edges inferred by similarity
    \item \textbf{Verbal critique}: Modifies node embedding position in semantic space
\end{itemize}

\textbf{Key Insight}: Hodge decomposition does \emph{not} obviate preference data---it \emph{processes} it. But it reduces collection burden via:
\begin{itemize}[leftmargin=*]
    \item \textbf{Transitive closure}: $A > B$ and $B > C$ implies $A > C$ via gradient
    \item \textbf{Automatic cycle detection}: $A > B > C > A$ captured in $\textcolor{harmcolor}{\mathbf{h}}$
    \item \textbf{Noise robustness}: Random errors appear in $\textcolor{curlcolor}{\text{curl}}$, not gradient
\end{itemize}

\subsection{Schwarzschild Detection: Identifying Black Holes}

How do we identify catastrophic regions from preference data alone? Ordinal preferences (``A > B > C'') provide only \emph{relative} information---they cannot distinguish ``mildly bad'' from ``catastrophically unsafe.''

\textbf{Three Complementary Methods}:

\begin{enumerate}[leftmargin=*]
    \item \textbf{Sink Detection}: Black holes act as flow sinks (high inflow, no outflow). Detect via divergence: $\text{div}(\nabla\phi)_i < -\tau$
    
    \item \textbf{Cliff Detection}: Catastrophic states exhibit steep gradients. Detect edges with $|Y_{ij}| > \tau_{cliff}$
    
    \item \textbf{Explicit Markers} (Recommended): Augment feedback with catastrophic flags (``A > B; C is \textbf{forbidden}'')
\end{enumerate}

\textbf{Schwarzschild Radius}: For detected black hole $b$, define event horizon:
\begin{equation}
    r_s(b) = \alpha \cdot \max_{i \to b} |Y_{ib}| + \beta \cdot \text{negative\_density}(b)
\end{equation}

\textbf{Limitation}: Pure ordinal feedback cannot distinguish ``bad'' from ``catastrophic.'' For safety-critical applications, explicit catastrophic labeling is strongly recommended.

\subsection{Topological Ranking via Hodge Gradients}

Instead of predicting a scalar reward, the \textbf{Hodge Critic} assigns a topological gradient $\textcolor{gradcolor}{\nabla\phi}$ to each action---the ``clean'' component after removing inconsistencies:

\begin{equation}
    \text{score}(a) = \langle \text{embed}(a), \textcolor{gradcolor}{\nabla\phi} \rangle
\end{equation}

Higher alignment = action moves along the geodesic toward genuinely preferred outcomes.

\begin{figure}[H]
    \centering
    \includegraphics[width=0.8\textwidth]{../../fig3_hodge_decomp.png}
    \caption{Hodge decomposition separates preference flow into $\textcolor{gradcolor}{\text{gradient}}$ (consistent), $\textcolor{curlcolor}{\text{curl}}$ (local cycles), and $\textcolor{harmcolor}{\text{harmonic}}$ (global cycles) components. Standard RLHF conflates all three; we extract only the gradient.}
    \label{fig:hodge_decomp}
\end{figure}

\section{Geometric Safety via Black Holes}

Traditional constrained RL treats safety as a cost limit $\mathbb{E}[C] \le d$, allowing catastrophic failures with low probability. We propose embedding safety directly into the \emph{geometry} of the state space.

\subsection{Semantic Failure Clusters}

A single catastrophic data point is insufficient to define a safety boundary. Dangerous behaviors (deception, sandbagging) cluster semantically. We aggregate black hole states into clusters $\mathcal{C} = \{C_1, \dots, C_m\}$ with centroids $\mu_j$ and event horizon radii $R_j$, calculated to enclose all member event horizons:

\begin{equation}
    R_j = \max_{x \in C_j} \left( \|x - \mu_j\| + r_s(x) \right) + \epsilon
\end{equation}

\subsection{Generalized Safety Metric}

We define the Riemannian metric $g_{ij}(s)$ using these clusters:
\begin{equation}
g_{ij}(s) = \left( 1 + \sum_{C_j \in \mathcal{C}} \frac{\alpha \cdot R_j}{d(s, \mu_j)^k} \right) \delta_{ij}
\end{equation}

Where $R_j$ ensures larger clusters exert repulsion from further away, creating spherical event horizons around entire failure modes.

\textbf{Theorem} (Geodesic Avoidance): For $k \ge 2$, the geodesic distance from any safe state to any point inside the event horizon ($d(s, \mu_j) < R_j$) is \emph{infinite}. Any policy minimizing geodesic path length avoids the entire class of behaviors represented by $C_j$ with probability 1.

\subsection{Intuition: The Supermassive Black Hole}

Why cluster failures? Consider the difference between ``Point-Based Safety'' (Sniper) and ``Cluster-Based Safety'' (Shotgun).

\begin{itemize}[leftmargin=*]
    \item \textbf{Point-Based (Sniper)}: Marking specific bad responses creates ``pinhole'' singularities. Adversarial policies can ``thread the needle''---finding paths $\epsilon$-close to catastrophe but technically unmarked.
    \item \textbf{Cluster-Based (Shotgun)}: Real failures (e.g., deception) are semantic regions. Aggregating points into a \textbf{Supermassive Black Hole} creates a continuous event horizon engulfing the entire neighborhood. Even novel variations of deceptive behavior fall within the cluster's radius.
\end{itemize}

\begin{figure}[H]
    \centering
    \includegraphics[width=0.8\textwidth]{../../fig2_geometric_safety.png}
    \caption{Geometric safety: ``Black hole'' regions have infinite geodesic distance, creating hard safety guarantees without explicit constraints. Unlike CPO's soft penalties, this barrier cannot be overcome by high reward.}
    \label{fig:geometric_safety}
\end{figure}

\section{Sheaf-Geodesic Policy Optimization (SGPO)}

\subsection{Algorithm Overview}

SGPO combines Hodge-decomposed rewards with geometric safety:
\begin{enumerate}[leftmargin=*]
    \item \textbf{Collect trajectories} from current policy $\pi_{\theta}$
    \item \textbf{Compute Hodge decomposition} of preference feedback
    \item \textbf{Extract gradient component} $\textcolor{gradcolor}{\nabla\phi}$ as the ``clean'' reward signal
    \item \textbf{Update safety metric} based on identified black holes
    \item \textbf{Optimize policy} to follow $\textcolor{gradcolor}{\nabla\phi}$ while respecting geodesic constraints
\end{enumerate}

\subsection{Geodesic DPO Loss}

Instead of standard DPO, we use:
\begin{equation}
    \mathcal{L}_{\text{GeoDPO}} = -\log \sigma\left(\beta \cdot \langle \textcolor{gradcolor}{\nabla\phi}, \Delta e \rangle_g\right)
\end{equation}

Where $\langle \cdot, \cdot \rangle_g$ is the inner product under the safety metric $g$, automatically down-weighting updates that would approach black holes.

\section{Embedding Representations}

Our framework operates on \emph{any} embeddable state representation. While natural language is convenient for LLM-based agents, structured representations (JSON schemas, typed APIs, domain-specific languages) often provide \emph{more consistent} embeddings with fewer spurious correlations.

\begin{table}[H]
    \centering
    \begin{tabular}{lccc}
        \toprule
        \textbf{Representation} & \textbf{Embedding} & \textbf{Semantic} & \textbf{Consistency} \\
        & \textbf{Difficulty} & \textbf{Richness} & \textbf{(Lower $H^1$)} \\
        \midrule
        Raw pixels & Hard (CNN) & Low & Variable \\
        Structured JSON/APIs & Medium & Medium & \textbf{High} \\
        Natural language & Easy & High & Low \\
        Hybrid (NL + schema) & Medium & High & Medium \\
        \bottomrule
    \end{tabular}
    \caption{State representation tradeoffs. Structured formats reduce embedding inconsistencies ($H^1$) at the cost of flexibility.}
    \label{tab:representations}
\end{table}

\textbf{Key Insight}: The same semantic concept (e.g., ``user satisfaction'') may embed differently depending on phrasing, introducing \emph{artificial} $H^1$ inconsistency. This is precisely the problem Hodge decomposition addresses---separating genuine preference cycles from representation artifacts.

\section{Preliminary Results}

We evaluate STRS on synthetic benchmarks designed to isolate preference cycles and safety violations.

\subsection{Condorcet Cycle Detection}

On a circular state space with cyclic reward structure:
\begin{table}[H]
    \centering
    \begin{tabular}{lcc}
        \toprule
        \textbf{Method} & \textbf{Learned $H^1$} & \textbf{Ground Truth $H^1$} \\
        \midrule
        PPO (implicit) & N/A (exploits cycle) & 0.500 \\
        SGPO (ours) & 0.357 & 0.500 \\
        \bottomrule
    \end{tabular}
    \caption{SGPO recovers 71\% of the true cycle magnitude; PPO has no mechanism to detect cycles.}
\end{table}

\subsection{Safety Benchmark}

On a navigation task with deceptive high-reward trap:
\begin{table}[H]
    \centering
    \begin{tabular}{lccc}
        \toprule
        \textbf{Method} & \textbf{Mean Return} & \textbf{Violations} & \textbf{Reduction} \\
        \midrule
        PPO & -6.67 & 52 & --- \\
        CPO & -6.23 & 7 & 87\% \\
        SGPO (ours) & \textbf{+1.53} & \textbf{11} & \textbf{79\%} \\
        \bottomrule
    \end{tabular}
    \caption{SGPO achieves positive return while reducing violations by 79\% vs.\ PPO.}
\end{table}

\section{Conclusion}

We introduced \textbf{Sheaf-Theoretic Reward Spaces}, a framework that makes explicit the geometric structure implicit in existing RLHF methods. By decomposing preferences into gradient, curl, and harmonic components, we can:
\begin{enumerate}[leftmargin=*]
    \item \textbf{Detect} cyclic inconsistencies that cause reward hacking
    \item \textbf{Quantify} the degree of evaluator disagreement via $H^1$ cohomology
    \item \textbf{Enforce} hard safety guarantees via geometric barriers
\end{enumerate}

PPO's success can be understood as implicit local navigation of the reward manifold; CPO adds soft constraint boundaries. SGPO makes this geometry explicit, enabling principled handling of the topological defects that cause alignment failures.

\textbf{Future Work}:
\begin{itemize}[leftmargin=*]
    \item \textbf{Multi-Objective Extension}: The current framework assumes a single global preference ordering. For \emph{multi-objective} settings (Pareto frontiers), we need \textbf{vector-valued Hodge decomposition} where each objective dimension has its own gradient field. This connects to multi-task learning and constitutional AI.
    
    \item \textbf{Scalability}: Applying cohomology detection to frontier-scale RLHF datasets with millions of comparisons.
    
    \item \textbf{Learned Restriction Maps}: Currently, we use embedding similarity to infer edges. Learning task-specific restriction maps could improve context-dependent preference handling.
    
    \item \textbf{Schwarzschild Estimation}: Developing better heuristics for inferring black hole severity from ordinal data alone, potentially via uncertainty quantification or ensemble disagreement.
\end{itemize}

\bibliographystyle{plainnat}
\begin{thebibliography}{9}
\bibitem[Schulman et al.(2017)]{schulman2017proximal}
Schulman, J., Wolski, F., Dhariwal, P., Radford, A., \& Klimov, O. (2017).
\newblock Proximal Policy Optimization Algorithms.
\newblock \emph{arXiv:1707.06347}.

\bibitem[Achiam et al.(2017)]{achiam2017constrained}
Achiam, J., Held, D., Tamar, A., \& Abbeel, P. (2017).
\newblock Constrained Policy Optimization.
\newblock \emph{ICML}.

\bibitem[Condorcet(1785)]{condorcet1785}
Condorcet, M. (1785).
\newblock Essai sur l'application de l'analyse \`a la probabilit\'e des d\'ecisions.

\bibitem[Jiang et al.(2011)]{jiang2011statistical}
Jiang, X., Lim, L.-H., Yao, Y., \& Ye, Y. (2011).
\newblock Statistical Ranking and Combinatorial Hodge Theory.
\newblock \emph{Mathematical Programming}, 127(1), 203--244.

\bibitem[Ouyang et al.(2022)]{ouyang2022training}
Ouyang, L., et al. (2022).
\newblock Training language models to follow instructions with human feedback.
\newblock \emph{NeurIPS}.
\end{thebibliography}

\end{document}
